\section{Caso de estudio: nanoantenas}
\label{sec:nanoantenas}

A lo largo del trabajo se ha caracterizado la dispersión de luz por esferas
de oro aisladas y por agregados de dos y cuatro esferas en términos de dos
magnitudes: las eficiencias espectrales en campo lejano
($Q_\text{abs}$, $Q_\text{sca}$, $Q_\text{ext}$) y la intensificación de
campo cercano $|\mathbf{E}|^2$ en el entorno de las partículas. Estos
resultados pueden reinterpretarse dentro del marco de las
\emph{nanoantenas ópticas}. Esta sección presenta dicho concepto, lo
conecta con el formalismo desarrollado en la sección~\ref{sec:miecoeffs} y
reinterpreta los resultados de las secciones~\ref{sec:dimero}
y~\ref{sec:tetramero} desde este punto de vista.

% =========================================================================
\subsection{¿Qué es una nanoantena óptica?}
\label{subsec:que_es_nanoantena}
% =========================================================================

En radiofrecuencia y microondas, una antena es un dispositivo que convierte
la energía electromagnética que viaja confinada en un circuito o guía en
radiación que se propaga libremente, y viceversa~\cite{Novotny2011}. Su
función es servir de interfaz entre campos localizados y campos propagantes.
Una \emph{nanoantena óptica} desempeña un papel análogo a frecuencias
ópticas —concentrar la luz incidente en volúmenes inferiores a la
longitud de onda, generando una región de campo intenso o \emph{hot spot},
y, recíprocamente, transformar la energía de una fuente localizada (por
ejemplo, una molécula) en radiación al campo lejano—, aunque condicionado
por la respuesta dispersiva y disipativa de los metales a frecuencias
ópticas~\cite{Bharadwaj2009,Novotny2011}.

% La diferencia esencial respecto a las antenas convencionales es que, a
% frecuencias ópticas, los metales no pueden describirse como conductores
% perfectos. Cuando la luz incide sobre una partícula metálica pequeña, su
% campo eléctrico empuja de forma colectiva a los electrones de conducción, que
% oscilan de un lado a otro de la partícula. A ciertas frecuencias esta
% oscilación entra en resonancia: es lo que se conoce como \emph{plasmón de
% superficie localizado} (LSP), y es el responsable de que el campo se
% intensifique en las inmediaciones de la partícula. En el formalismo de Mie
% empleado en este trabajo, esta respuesta queda recogida en los coeficientes
% de dispersión $a_n$ (sección~\ref{sec:miecoeffs}): las resonancias
% plasmónicas aparecen como máximos de dichos coeficientes, siendo el término
% dipolar $a_1$ el dominante para las partículas más pequeñas.

A diferencia de las antenas convencionales, que a frecuencias de radio y
microondas pueden aproximarse como conductores perfectos, a frecuencias
ópticas los metales presentan una permitividad finita, compleja y dependiente
de la frecuencia. Cuando la luz incide sobre una partícula metálica pequeña,
su campo eléctrico empuja de forma colectiva a los electrones de conducción,
que oscilan de un lado a otro de la partícula. Esta oscilación entra en
resonancia cuando la permitividad compleja del metal alcanza ciertos valores,
que en metales nobles como el oro suelen darse en el visible o cerca del
visible: es lo que se conoce como \emph{plasmón de superficie localizado}
(LSP), y es el responsable de que el campo se intensifique en las
inmediaciones de la partícula. En el formalismo de Mie empleado en este
trabajo, esta respuesta queda recogida en los coeficientes de dispersión
$a_n$ (sección~\ref{sec:miecoeffs}): las resonancias plasmónicas aparecen como
máximos de dichos coeficientes, siendo el término dipolar $a_1$ el dominante
para las partículas más pequeñas.

Estos máximos resonantes están ligados a las frecuencias de operación de la
antena, mientras que el carácter de la resonancia (dipolar, cuadrupolar,
etc.) condiciona el patrón de radiación, es decir, la distribución angular de
la luz dispersada (véase la Figura~\ref{fig:S11_regimenes}). Cuando el
parámetro de tamaño es pequeño y domina el modo dipolar, una esfera metálica
puede interpretarse como una antena dipolar resonante elemental~\cite{Kuhn2006}.

Conviene distinguir dos funciones reconocibles en los
resultados de este trabajo~\cite{Kolwas2013}:

\begin{itemize}
  \item \textbf{Antena receptora.} La partícula transforma parte de la
        radiación incidente en energía localizada en su entorno inmediato,
        intensificando el campo cercano. En partículas pequeñas, este
        comportamiento suele estar acompañado por una contribución importante
        de la absorción, debido al predominio de las pérdidas disipativas.
  \item \textbf{Antena emisora.} La partícula puede actuar también como canal
        radiativo, reemitiendo energía hacia el campo lejano. Este carácter
        radiativo se manifiesta en una mayor contribución de la dispersión,
        especialmente al aumentar el tamaño de la partícula.
\end{itemize}

El predominio relativo de estos comportamientos depende del tamaño de la
partícula, de la longitud de onda, del medio circundante y, en agregados, de
la separación entre partículas. En particular, al aumentar el tamaño la
partícula tiende a reemitir la luz incidente en lugar de disiparla en forma
de calor, lo que permite interpretar el cambio relativo entre $Q_\text{abs}$
y $Q_\text{sca}$ observado al pasar de $a=25\,\text{nm}$ a $a=85\,\text{nm}$
en la sección~\ref{sec:dimero} como una transición desde una respuesta más
disipativa hacia otra de mayor carácter radiativo.

% =========================================================================
\subsection{La esfera y el dímero como nanoantena}
\label{subsec:esfera_dimero_antena}
% =========================================================================

La esfera individual constituye la nanoantena más simple. Su resonancia
plasmónica dipolar se asocia con una frecuencia de operación que puede
sintonizarse variando el radio: al aumentar el tamaño, la resonancia se
desplaza al rojo y crece la contribución de modos de orden superior
(cuadrupolar y siguientes), tal y como se observó en los espectros de
referencia de la sección~\ref{sec:dimero}. Sin embargo, en las
configuraciones estudiadas la esfera aislada produce una intensificación
moderada, del orden de diez veces la intensidad incidente, porque no presenta
un gap donde se concentre el campo por acoplamiento entre partículas.

El dímero introduce precisamente esa región. El gap desempeña un papel
análogo al punto de alimentación de una antena, al ser la región donde se
localiza con mayor intensidad el campo cercano por el acoplamiento de las dos
partículas~\cite{Yousif2013}. Una magnitud para cuantificar este
efecto es el factor de intensificación de campo en el gap,
$|\mathbf{E}|^2/|\mathbf{E}_0|^2$, equivalente al \emph{local field intensity
enhancement factor} (LFIEF) empleado en la literatura~\cite{Yousif2013}.
La iluminación empleada en este trabajo es una onda monocromática no
polarizada. Aun así, como el dímero está orientado con su eje a lo largo de
la dirección $x$ —contenida en el plano de polarización—, de las dos
componentes ortogonales del campo es la paralela al eje la que excita el modo
de acoplamiento entre las esferas y concentra el campo en el gap. Al estrechar
el gap aumenta este acoplamiento de campo cercano, lo que se traduce en una
mayor intensificación en el hueco: en la sección~\ref{sec:dimero}, el dímero
heterogéneo pasó de $21$ a $186\,|\mathbf{E}_0|^2$ al reducir $d$ de $50$ a
$5\,\text{nm}$. El efecto sobre la posición espectral de la resonancia, en
cambio, depende del tamaño: es prácticamente nulo en el dímero de
$25\,\text{nm}$, mientras que en el de $85\,\text{nm}$ la resonancia de
$Q_\text{sca}$ se desplaza al rojo en más de $170\,\text{nm}$ al cerrar el
gap. Esta dependencia con el tamaño es coherente con las simulaciones de
dímeros de oro mediante el método de la T-matrix~\cite{Yousif2013}.

Conviene señalar que los resultados de referencia más próximos a este trabajo
emplean la misma herramienta numérica: las cadenas de esferas de oro de
Harris \textit{et al.} se calcularon con el código de la T-matrix de
Mackowski y Mishchenko~\cite{Harris2009}, 
el mismo formalismo descrito en la
sección~\ref{sec:fundamento multiple} e implementado en MSTM. Allí, empleando
polarización paralela al eje de la cadena, se observa que en el régimen de
campo cercano fuertemente acoplado la resonancia longitudinal se desplaza al
rojo de forma aproximadamente proporcional a $1/d$ al reducir la
separación~\cite{Harris2009}, en línea con el corrimiento al rojo que
observamos al cerrar el gap en el dímero de $85\,\text{nm}$.

% =========================================================================
\subsection{Reinterpretación de los resultados como nanoantenas}
\label{subsec:reinterpretacion}
% =========================================================================

La Tabla~\ref{tab:antena_merito} resume las configuraciones representativas
ordenadas por su intensificación máxima de campo cercano. Esta intensificación
es solo uno de los muchos criterios con los que puede caracterizarse una
nanoantena: también pueden considerarse su eficiencia radiativa, su
directividad o su sección eficaz de dispersión, entre otros. El canal de
extinción dominante en cada caso —absorción o dispersión— se recoge en las
Tablas~\ref{tab:dimero_albedo} y~\ref{tab:tetramero_albedo}: las
configuraciones con esferas de $25\,\text{nm}$ permanecen dominadas por la
absorción, mientras que las que incorporan esferas de $85\,\text{nm}$
presentan un mayor peso del canal radiativo (dispersión).
 
\begin{table}[H]
  \centering
  \begin{tabular}{lccc}
    \hline
    Configuración ($d=5\,\text{nm}$ en agregados) & $\lambda_\text{res}$ &
    $|\mathbf{E}|^2_\text{máx}/|\mathbf{E}_0|^2$ & Canal de extinción \\
    \hline
    Esfera individual, $a=25\,\text{nm}$        & 520\,nm & $11.9$  & Absorción \\
    Esfera individual, $a=85\,\text{nm}$        & 570\,nm & $12.8$  & Dispersión \\
    Dímero homogéneo, $a=25\,\text{nm}$         & 530\,nm & $147$   & Absorción \\
    Dímero heterogéneo, $a=25/85\,\text{nm}$    & 590\,nm & $186$   & Dispersión \\
    Tetrámero homogéneo, $a=25\,\text{nm}$      & 550\,nm & $431$   & Absorción \\
    Tetrámero heterogéneo, $a=85/25/25/85$      & 620\,nm & $566$   & Dispersión \\
    \hline
  \end{tabular}
  \caption{Intensificación máxima de campo cercano de las configuraciones
           más representativas, interpretadas como nanoantenas. El canal de
           extinción indica si en la longitud de onda resonante domina la
           absorción o la dispersión
           (Tablas~\ref{tab:dimero_albedo} y~\ref{tab:tetramero_albedo}). El
           paso de la esfera al dímero y de este al tetrámero incrementa la
           concentración de campo.}
  \label{tab:antena_merito}
\end{table}

Dos tendencias destacan al interpretar los resultados como una antena. La primera
es el papel del gap como región de máxima concentración de campo: a igualdad
de tamaño y número de esferas, estrechar el hueco se asocia con una mayor
intensificación, al aumentar el acoplamiento de campo cercano entre
partículas. La segunda es el efecto de añadir elementos a la cadena. Al pasar
del dímero al tetrámero homogéneo de $25\,\text{nm}$, la intensificación en el
gap central crece de $147$ a $431\,|\mathbf{E}_0|^2$ y la resonancia se
desplaza al rojo, lo que indica que, para el modo longitudinal estudiado, las
partículas adicionales contribuyen al acoplamiento colectivo de la cadena.
Este comportamiento es coherente con el observado experimentalmente en
cadenas de oro de uno a seis elementos, donde el desplazamiento al rojo crece
con la longitud de la cadena y satura en torno a 10--12
partículas~\cite{Barrow2011}, y con los cálculos de T-matrix, según los cuales
—para separaciones pequeñas (gap $\lesssim 5\,\text{nm}$)— la resonancia
longitudinal se aproxima a su valor de cadena infinita ya hacia las $\sim\!10$
partículas~\cite{Harris2009}. El tetrámero estudiado aquí puede entenderse,
por tanto, como un segmento corto de esa serie, en la zona donde el
desplazamiento todavía crece al añadir partículas.

Desde esta misma perspectiva, dos aspectos merecen destacarse. El primero es que la
asimetría geométrica actúa como un grado de libertad de diseño: combinar
esferas de distinto tamaño no solo aumenta la intensificación, sino que
permite fijar la posición del hot spot dentro del agregado, en el lado de la
esfera menor. El segundo es que esta interpretación deja de ser aplicable al aumentar
suficientemente el tamaño: la esfera de $a=600\,\text{nm}$, ya en pleno
régimen multipolar (sección~\ref{sec:dimero}), no presenta una resonancia
localizada dominante y no concentra el campo de forma eficiente. El tamaño a
partir del cual se produce esta transición se sitúa por debajo de ese valor,
aunque no se ha acotado en este trabajo.

La idea de combinar esferas de distinto tamaño para concentrar el campo y
situar el \emph{hot spot} aparece ya en el concepto de \emph{nanolente
plasmónica}, una cadena de nanoesferas metálicas de tamaño decreciente que
focaliza la luz en el gap entre las más pequeñas~\cite{Li2003}. En las
versiones simétricas de esas cadenas, el campo se concentra en el gap
central. La configuración $85$--$25$--$25$--$85$ estudiada aquí sigue esa
misma idea. La diferencia es que aquel trabajo usa plata en régimen
cuasiestático, mientras que aquí se emplean esferas de oro mayores y
simulación de onda completa (MSTM). Aun así, se reproduce el mismo 
comportamiento cualitativo: la
mayor intensificación se da en el gap central, entre las esferas pequeñas.

% =========================================================================
\subsection{Aplicaciones}
\label{subsec:aplicaciones_antena}
% =========================================================================
 
El interés práctico de las nanoantenas plasmónicas reside en que los
\emph{hot spots} generan campos locales mucho más intensos que el campo
incidente. Una molécula situada en el gap entre partículas experimenta,
por tanto, una interacción reforzada con la luz, lo que puede aumentar de
forma significativa la señal óptica observada. Tres ejemplos ilustran la
utilidad de este mecanismo.
 
\begin{itemize}
  \item \textbf{Espectroscopía Raman amplificada por superficie (SERS).}
        Cuando la luz ilumina una molécula, una pequeña fracción se dispersa
        con una frecuencia distinta de la incidente, determinada por los
        modos vibracionales de la molécula. Este espectro Raman actúa como
        una \emph{huella dactilar} que permite identificarla. El problema es
        que la señal Raman espontánea es muy débil, por lo que normalmente
        se requiere una cantidad apreciable de moléculas para medirla. Si la
        molécula se sitúa en un \emph{hot spot}, el campo concentrado puede
        amplificar notablemente la señal, llegando en condiciones optimizadas
        a la detección de cantidades extremadamente pequeñas de analito e
        incluso de moléculas individuales. Este principio se emplea, por
        ejemplo, en la detección de trazas químicas, contaminantes o residuos
        de pesticidas en muestras muy reducidas~\cite{Pilot2019}.

  \item \textbf{Realce de fluorescencia.} Muchas moléculas absorben luz y la
        reemiten con otro color; este fenómeno, la fluorescencia, se usa en
        biología para «iluminar» estructuras concretas de una célula al
        microscopio, marcándolas con moléculas fluorescentes. Cuanto más
        brilla el marcador, más fácil es verlo. Una nanoantena situada al lado
        hace que la molécula emita más luz 
        actuando como intermediaria entre el emisor localizado y el campo lejano, 
        de modo que un marcador demasiado tenue puede
        volverse visible; esto se ha comprobado acercando una nanopartícula de
        oro a una sola molécula. Eso sí, si se acerca demasiado al metal la
        energía se disipa en él en forma de calor y la molécula se apaga en
        lugar de brillar más, por lo que existe una distancia
        óptima~\cite{Kuhn2006}.

  \item \textbf{Medida de distancias.} La longitud de onda de la resonancia
        de un dímero depende de la separación entre sus partículas: al
        reducir el gap, el acoplamiento de campo cercano aumenta y la
        resonancia se desplaza hacia el rojo. Esta dependencia
        puede utilizarse de forma inversa, deduciendo la distancia entre
        partículas a partir del desplazamiento espectral, lo que da lugar a
        las denominadas \emph{reglas plasmónicas}. Así, al unir nanopartículas
        metálicas a distintas partes de una biomolécula, los cambios en el
        espectro pueden informar sobre variaciones en su forma, como
        procesos de plegamiento o estiramiento, en escalas por debajo del
        límite de resolución de un microscopio óptico convencional~\cite{Sonnichsen2005}.
\end{itemize}
 
En conjunto, los sistemas simulados en este trabajo —de la esfera aislada al
tetrámero— reproducen los ingredientes físicos en los que se apoyan estas
aplicaciones: la concentración del campo, el desplazamiento de la resonancia
y el control de la respuesta mediante el tamaño, el gap y la forma del
agregado.